%% kapa: exam=1 round=21 year=2010 period=2 form=A subject=accounting subject_ko=회계학 seq=26 number=26
%% source: Q-Net 감정평가사 기출문제 (cst00309) artlSeq=5024181 file=extracted/1008/5024181/2010년 제21회 감정평가사 1차시험 기출문제/2교시.hwp
%% conversion: kordoc 4.24.51 -> tools/gen1.py (auto); review=proofread-gpt-sol; flags=none
\begin{question}{26}
한국채택국제회계기준에서 연결재무제표와 별도재무제표의 용어에 관한 설명으로 옳지 않은 것은?
\begin{choices}
  \item 종속기업은 다른 기업의 지배를 받고 있는 기업으로서 파트너십과 같은 법인격 없는 실체를 제외한다.
  \item 지배기업은 하나 이상의 종속기업을 가지고 있는 기업이다.
  \item 지배력은 경제활동에서 효익을 얻기 위하여 재무정책과 영업정책을 결정할 수 있는 능력이다.
  \item 비지배지분은 종속기업의 지분 중 지배기업에게 직접으로 또는 간접으로 귀속되지 않는 지분이다.
  \item 별도재무제표는 지배기업, 관계기업의 투자자 또는 공동지배기업의 참여자가 투자자산을 피투자자의 보고된 성과와 순자산에 근거하지 않고 직접적인 지분투자에 근거한 회계처리로 표시한 재무제표이다.
\end{choices}
\end{question}
\answer{1}
